% Pursuing Stacks (cotes 134-2 to 134-8): archive cote 134-3, site page 23 (PDF page 24), typescript page 23.
% typescript. Draft compared: ps2.tex:1055-1122.
% AI-transcribed by Claude Opus 5.5 (claude-opus-5-5, adaptive thinking, effort high), Message Batches API, batch ps1 (msgbatch_01Rv5hCaNTtUt4JHx1yqoHPL), 2026-11-12; not yet checked by eye.
\scanpage{134-3/23}
\mspage{23}
a model category in Quillen's sense,  the situation with the inclusion
functor
\[ N:(Cat)\longrightarrow (Ss sets) \]
and the left adjoint functor N'  is by no means the one of Quillen\hand{'}s
comparison theorem,  where the two functors play\xout{ed} mutually dual roles
and both induce equivalences on the corresponding localized homotopical
categories.  Here \uline{only} N,  \uline{not} N',  \uline{induces such equivalence}.  This is
analoguous to the situation,  already note\over{\unsure{s}}{d} before,  of the inclusion
functors
\[ (Preord)\hookrightarrow (Cat) \qquad or \qquad (Ord)\hookrightarrow (Cat) \quad , \]
which \inserted{[presumably]} by localization induce equivalences on the associated homotopy
categories,  but the left adjoint functors do not share this property.

The topology $T_\infty$  just considered on (Cat) as a ``suitable'' topology
for describing homotopy types of objects of (Cat),  was\hand{|}of course direct-%
ly inspired by the semi-simplicial approach to homotopy types,  via
simplicial complexes  namely via the two associated inclusion functors
\[ \hand{\Delta} \hookrightarrow (Cat)\hookrightarrow \hand{\Delta}^{\wedge} , \]
the second functor associating to every category A the ``complex'' of it's
``simplicial diagrams''  $a_o\longrightarrow a_1\longrightarrow \;\ldots\longrightarrow a_n$ .  If we had been working
\xout{rather} with  cubical complexes rather than ss ones for describing homoto
-py types,  this would give rise similarly to two functors
\[ \hand{\square} \hookrightarrow (Cat)\hookrightarrow \hand{\square}^{\wedge} , \]
where $\hand{\square}$ is the category of``standar\over{t}{d} cubes'' and face and degeneracy
maps between them,  and where the second\hand{|}inclusion associates to every
category A the cubical complex of it's  ``cubical diagrams'',  namely \inserted[typed]{commutative}
diagrams in A modelled after the diagram types $\hand{\square}_n\hand{[}1\hand{]}$ ,  the 1-skeleton
of the standart n-cubes $\hand{\square}_n$ , with suitable orientations on it's edges
(indicative of the direction of corresponding arrows in A).  The natural
idea \over{\xout{here}}{now} would be to endow (Cat) with the \over{category}{topology}, $T'_\infty$ say,  induced
by $\hand{\square}^{\wedge}$,  namely call a family (*)  covering iff for every $n\hand{\in} \underline{N}$,  the
corresponding family  of \xout{set} maps of sets
\[ Cub_n(A_i)\longrightarrow \hand{C}ub_n(A) \]
is epimorphic.   This \xout{category} topology appears to be \xout{finer}\inserted[typed]{coarser} than $T_\infty$ \xout{,}
\inserted[typed]{(i.e. there are less covering families),}
and the comparison lemma gives now the equivalence
\[ ((Cat),T'_\infty)\hand{^{\sim}} \hand{\simeq} \hand{\square}^{\wedge} , \]
which shows that definitely the topology is \uline{strictly} coarser than $T_\infty$ ,
as it gives rise to a non-equivalent topos,  $\hand{\square}^{\wedge}$ instead of $\hand{\Delta}^{\wedge}$.

T\hand{h}ese reflections convince me  1)  there are indeed topologies
on (Cat),  suit\unsure{able for} describing the natural homotopy types of ob-%
jects of (Cat) namely of categories,  and  2)  that there is definitely
